% Options for packages loaded elsewhere
\PassOptionsToPackage{unicode}{hyperref}
\PassOptionsToPackage{hyphens}{url}
\documentclass[
  10pt,
  fontset=fandol]{ctexart}
\usepackage{xcolor}
\usepackage[margin=16mm]{geometry}
\usepackage{amsmath,amssymb}
\setcounter{secnumdepth}{-\maxdimen} % remove section numbering
\usepackage{iftex}
\ifPDFTeX
  \usepackage[T1]{fontenc}
  \usepackage[utf8]{inputenc}
  \usepackage{textcomp} % provide euro and other symbols
\else % if luatex or xetex
  \usepackage{unicode-math} % this also loads fontspec
  \defaultfontfeatures{Scale=MatchLowercase}
  \defaultfontfeatures[\rmfamily]{Ligatures=TeX,Scale=1}
\fi
\usepackage{lmodern}
\ifPDFTeX\else
  % xetex/luatex font selection
\fi
% Use upquote if available, for straight quotes in verbatim environments
\IfFileExists{upquote.sty}{\usepackage{upquote}}{}
\IfFileExists{microtype.sty}{% use microtype if available
  \usepackage[]{microtype}
  \UseMicrotypeSet[protrusion]{basicmath} % disable protrusion for tt fonts
}{}
\makeatletter
\@ifundefined{KOMAClassName}{% if non-KOMA class
  \IfFileExists{parskip.sty}{%
    \usepackage{parskip}
  }{% else
    \setlength{\parindent}{0pt}
    \setlength{\parskip}{6pt plus 2pt minus 1pt}}
}{% if KOMA class
  \KOMAoptions{parskip=half}}
\makeatother
\setlength{\emergencystretch}{3em} % prevent overfull lines
\providecommand{\tightlist}{%
  \setlength{\itemsep}{0pt}\setlength{\parskip}{0pt}}
\usepackage{amsmath,amssymb,mathtools}
\setcounter{secnumdepth}{0}
\setlength{\emergencystretch}{2em}
\allowdisplaybreaks
\usepackage{bookmark}
\IfFileExists{xurl.sty}{\usepackage{xurl}}{} % add URL line breaks if available
\urlstyle{same}
\hypersetup{
  pdftitle={ch07b 转录语法校验稿},
  hidelinks,
  pdfcreator={LaTeX via pandoc}}

\title{ch07b 转录语法校验稿}
\author{}
\date{}

\begin{document}
\maketitle

\clearpage

原书 PDF 196；印刷页 183。

\href{/Users/xiangjunfeng/Desktop/数值计算专用/numerical-analysis/book/source-images/pdf-196.jpeg}{原页图像}

（续 7.4.2 平方根法。）

求解 \(Ax=b\)，即求解如下两个三角方程组：

（1）\(Ly=b\)，求 \(y\)；（2）\(L^{\mathrm T}x=y\)，求 \(x\)。

\[
\begin{aligned}
\text{步 3}\quad &y_i=\left(b_i-\sum_{k=1}^{i-1}l_{ik}y_k\right)/l_{ii}\quad(i=1,2,\cdots,n).\\
\text{步 4}\quad &x_i=\left(b_i-\sum_{k=i+1}^{n}l_{ki}x_k\right)/l_{ii}\quad(i=n,n-1,\cdots,1).
\end{aligned}
\tag{7.4.8}
\]

由式（7.4.7）知
\(a_{jj}=\sum_{k=1}^{j}l_{jk}^{2}\)（\(j=1,2,\cdots,n\)），所以

\[
l_{jk}^{2}\leqslant a_{jj}\leqslant\max_{1\leqslant j\leqslant n}\{a_{jj}\},
\]

\[
\max_{j,k}\{l_{jk}^{2}\}\leqslant\max_{1\leqslant j\leqslant n}\{a_{jj}\}.
\]

上面分析说明，分解过程中元素 \(l_{jk}\) 的数量级不会增长且对角元素
\(l_{jj}\) 恒为正数。于是不选主元素的平方根法是一个数值稳定的方法。

当求出 \(L\) 的第 \(j\) 列元素时，\(L^{\mathrm T}\) 的第 \(j\)
行元素亦就算出，所以平方根法约需 \(n^3/6\) 次乘除法运算，大约为一般直接
LU 分解法计算量的一半。

由于 \(A\) 为对称阵，因此在用计算机计算时只需存储 \(A\)
的下三角部分，共需要存储 \(n(n+1)/2\) 个元素，可用一维数组存放，即
\(\{a_{11},a_{21},a_{22},\cdots,a_{n1},a_{n2},\cdots,a_{nn}\}\)。

矩阵元素 \(a_{ij}\) 存放于一维数组的第
\(\left(i\cdot\frac{i-1}{2}+j\right)\) 个元素，\(L\) 的元素存放在 \(A\)
的相应位置。

由式（7.4.7）可看出，用平方根法解对称正定方程组时，计算 \(L\) 的元素
\(l_{ii}\) 需要用到开方运算，为了避免开方运算，下面用定理 7.7 的分解式
\(A=LDL^{\mathrm T}\)，即

\[
A=
\begin{pmatrix}
1&&&&\\
l_{21}&1&&&\\
l_{31}&l_{32}&1&&\\
\vdots&\vdots&\ddots&\ddots&\\
l_{n1}&l_{n2}&\cdots&l_{n,n-1}&1
\end{pmatrix}
\begin{pmatrix}
d_1&&&&\\
&d_2&&&\\
&&\ddots&&\\
&&&\ddots&\\
&&&&d_n
\end{pmatrix}
\begin{pmatrix}
1&l_{21}&l_{31}&\cdots&l_{n1}\\
&1&l_{32}&\cdots&l_{n2}\\
&&\ddots&\ddots&\vdots\\
&&&\ddots&l_{n-1,n-1}\\
&&&&1
\end{pmatrix},
\]

按行计算 \(L\) 的元素
\(l_{ij}\)（\(j=1,2,\cdots,i-1\)）。由矩阵乘法，并注意
\(l_{jj}=1\)，\(l_{jk}=0\)（\(j<k\)），得

\[
a_{ij}=\sum_{k=1}^{n}(LD)_{ik}(L^{\mathrm T})_{kj}
=\sum_{k=1}^{n}l_{ik}d_kl_{jk}
=\sum_{k=1}^{j-1}l_{ik}d_kl_{jk}+l_{ij}d_jl_{jj}.
\]

于是得到以下计算 \(L\) 的元素及 \(D\) 的对角元素公式：

对于 \(i=1,2,\cdots,n\)，

\[
\begin{aligned}
\text{步 1}\quad &l_{ij}=\left(a_{ij}-\sum_{k=1}^{j-1}l_{ik}d_kl_{jk}\right)/d_j\quad(j=1,2,\cdots,i-1).\\
\text{步 2}\quad &d_i=a_{ii}-\sum_{k=1}^{i-1}l_{ik}^{2}d_k.
\end{aligned}
\tag{7.4.9}
\]

\begin{quote}
校注（agent 补充，原书疑误）：式（7.4.8）步 4 原图确印
\(b_i\)，这里忠实保留。由本页所列 \(L^{\mathrm T}x=y\)，回代右端应使用
\(y_i\)；否则与其前方三角方程组不一致。\href{/Users/xiangjunfeng/Desktop/数值计算专用/numerical-analysis/book/staging/ch07b/assets/196-step4.png}{原式局部图}。

校注（agent 补充，原书疑误）：\(LDL^{\mathrm T}\)
展开中最右矩阵次末行末项原图印作 \(l_{n-1,n-1}\)，这里保留原文；按左边
\(L\) 的转置，该位置应为
\(l_{n,n-1}\)。\href{/Users/xiangjunfeng/Desktop/数值计算专用/numerical-analysis/book/staging/ch07b/assets/196-ldlt.png}{原矩阵局部图}。
\end{quote}

\clearpage

原书 PDF 197；印刷页 184。

\href{/Users/xiangjunfeng/Desktop/数值计算专用/numerical-analysis/book/source-images/pdf-197.jpeg}{原页图像}

（续 7.4.2 平方根法。）

为避免重复计算，引进 \(t_{ij}=l_{ij}d_j\)，由式（7.4.9）得到以下按行计算
\(L,T\) 元素的公式：

对于 \(i=1,2,\cdots,n\)，

\[
\begin{aligned}
\text{步 1}\quad &t_{ij}=a_{ij}-\sum_{k=1}^{j-1}t_{ik}l_{jk}\quad(j=1,2,\cdots,i-1).\\
\text{步 2}\quad &l_{ij}=t_{ij}/d_j\quad(j=1,2,\cdots,i-1).\\
\text{步 3}\quad &d_i=a_{ii}-\sum_{k=1}^{i-1}t_{ik}l_{ik}.
\end{aligned}
\tag{7.4.10}
\]

计算出 \(T=LD\) 的第 \(i\) 行元素
\(t_{ij}\)（\(j=1,2,\cdots,i-1\)）后，存放在 \(A\) 的第 \(i\)
行相应位置，然后再计算 \(L\) 的第 \(i\) 行元素，存放在 \(A\) 的第 \(i\)
行，\(D\) 的对角元素存放在 \(A\) 的相应位置。例如，

\[
A=\begin{pmatrix}
a_{11}&a_{21}&a_{31}&a_{41}\\
a_{21}&a_{22}&a_{32}&a_{42}\\
a_{31}&a_{32}&a_{33}&a_{43}\\
a_{41}&a_{42}&a_{43}&a_{44}
\end{pmatrix}
\longrightarrow
\begin{pmatrix}
d_1&&&\\
l_{21}&d_2&&\\
l_{31}&l_{32}&d_3&\\
t_{41}&t_{42}&t_{43}&t_{44}
\end{pmatrix}
\longrightarrow
\begin{pmatrix}
d_1&&&\\
l_{21}&d_2&&\\
l_{31}&l_{32}&d_3&\\
l_{41}&l_{42}&l_{43}&d_4
\end{pmatrix}.
\]

对称正定矩阵 \(A\) 按 \(LDL^{\mathrm T}\) 分解和按 \(LL^{\mathrm T}\)
分解计算量差不多，但 \(LDL^{\mathrm T}\) 分解不需要开方计算。求解
\(Ly=b,DL^{\mathrm T}x=y\) 的计算公式分别为步 4、步 5 所述公式。

\[
\begin{aligned}
\text{步 4}\quad &\begin{cases}
y_1=b_1,\\
y_i=b_i-\displaystyle\sum_{k=1}^{i-1}l_{ik}y_h\quad(i=2,\cdots,n).
\end{cases}\\
\text{步 5}\quad &\begin{cases}
x_n=y_n/d_n,\\
x_i=y_i/d_i-\displaystyle\sum_{k=i+1}^{n}l_{ki}x_k\quad(i=n-1,\cdots,2,1).
\end{cases}
\end{aligned}
\tag{7.4.11}
\]

计算公式（7.4.10）、（7.4.11）称为\textbf{改进的平方根法}。

\subsubsection{7.4.3 追赶法}\label{ux8ffdux8d76ux6cd5}

在一些实际问题中，例如解常微分方程边值问题，解热传导方程以及船体数学放样中建立三次样条函数等中，都会要求解系数矩阵为对角占优的三对角方程组

\[
\begin{pmatrix}
b_1&c_1&&&&&\\
a_2&b_2&c_2&&&&\\
&\ddots&\ddots&\ddots&&&\\
&&a_i&b_i&c_i&&\\
&&&\ddots&\ddots&\ddots&\\
&&&&a_{n-1}&b_{n-1}&c_{n-1}\\
&&&&&a_n&b_n
\end{pmatrix}
\begin{pmatrix}x_1\\x_2\\\vdots\\x_i\\\vdots\\x_{n-1}\\x_n\end{pmatrix}
=
\begin{pmatrix}f_1\\f_2\\\vdots\\f_i\\\vdots\\f_{n-1}\\f_n\end{pmatrix},
\]

（方程组及说明续下页。）

\begin{quote}
校注（agent 补充，原书疑误）：式（7.4.11）步 4 的求和项原图印作
\(l_{ik}y_h\)，这里保留；求和指标为 \(k\) 且来源为 \(Ly=b\)，相应项应为
\(l_{ik}y_k\)。\href{/Users/xiangjunfeng/Desktop/数值计算专用/numerical-analysis/book/staging/ch07b/assets/197-solves.png}{原式局部图}。
\end{quote}

\clearpage

原书 PDF 198；印刷页 185。

\href{/Users/xiangjunfeng/Desktop/数值计算专用/numerical-analysis/book/source-images/pdf-198.jpeg}{原页图像}

（续 7.4.3 追赶法。）

简记作

\[
Ax=f,
\tag{7.4.12}
\]

其中 \(A\) 满足下列对角占优条件：

（1）\(|b_1|>|c_1|>0\)；

（2）\(|b_i|\geqslant|a_i|+|c_i|\)，\(a_i,c_i\neq0\)（\(i=2,3,\cdots,n-1\)）；

（3）\(|b_n|>|a_n|>0\)。

下面利用矩阵的直接三角分解法来推导解三对角方程组（7.4.12）的计算公式。由系数阵
\(A\) 的特点，可以将 \(A\) 分解为两个三角阵的乘积，即

\[
A=LU
\]

其中 \(L\) 为下三角阵，\(U\)
为单位上三角阵。下面说明这种分解是可能的。设

\[
A=\begin{pmatrix}
b_1&c_1&&&\\
a_2&b_2&c_2&&\\
&\ddots&\ddots&\ddots&\\
&&a_{n-1}&b_{n-1}&c_{n-1}\\
&&&a_n&b_n
\end{pmatrix}
=\begin{pmatrix}
\alpha_1&&&&\\
\gamma_2&\alpha_2&&&\\
&\gamma_3&\alpha_3&&\\
&&\ddots&\ddots&\\
&&&\gamma_n&\alpha_n
\end{pmatrix}
\begin{pmatrix}
1&\beta_1&&&\\
&1&\beta_2&&\\
&&\ddots&\ddots&\\
&&&1&\beta_{n-1}\\
&&&&1
\end{pmatrix},
\tag{7.4.13}
\]

其中 \(\alpha_i,\beta_i,\gamma_i\) 为待定系数。比较式（7.4.13）两端即得

\[
\left.\begin{aligned}
b_1&=\alpha_1,\quad c_1=\alpha_1\beta_1,\\
a_i&=\gamma_i,\quad b_i=\gamma_i\beta_{i-1}+\alpha_i\quad(i=2,\cdots,n),\\
c_i&=\alpha_i\beta_i\quad(i=2,3,\cdots,n-1).
\end{aligned}\right\}
\tag{7.4.14}
\]

由 \(\alpha_1=b_1\neq0\)，\(|b_1|>|c_1|>0\)，\(\beta_1=c_1/b_1\)，得
\(0<|\beta_1|<1\)。下面用归纳法证明

\[
|\alpha_i|>|c_i|\neq0\quad(i=1,2,\cdots,n-1),
\tag{7.4.15}
\]

即

\[
0<|\beta_i|<1,
\]

从而由式（7.4.14）可求出 \(\beta_i\)。

式（7.4.15）对 \(i=1\) 是成立的。现设式（7.4.15）对 \(i-1\) 成立，求证对
\(i\) 亦成立。

由归纳法假设 \(0<|\beta_{i-1}|<1\)，又由式（7.4.15）及 \(A\)
的假设条件，有

\[
|\alpha_i|=|b_i-a_i\beta_{i-1}|\geqslant|b_i|-|a_i\beta_{i-1}|>|b_i|-|a_i|\geqslant|c_i|\neq0,
\]

也就是 \(0<|\beta_i|<1\)。由式（7.4.14）得到

\[
\alpha_i=b_i-a_i\beta_{i-1}\quad(i=2,3,\cdots,n),\quad
\beta_i=c_i/(b_i-a_i\beta_{i-1})\quad(i=2,3,\cdots,n-1).
\]

这就是说，由 \(A\) 的假设条件完全确定了
\(\{\alpha_i\},\{\beta_i\},\{\gamma_i\}\)，实现了 \(A\) 的 LU 分解。

求解 \(Ax=f\) 等价于解两个三角方程组 \(Ly=f\) 与 \(Ux=y\)，先后求 \(y\)
与 \(x\)，从而得到以下解三对角方程组的\textbf{追赶法公式}：

\textbf{步 1} 计算 \(\{\beta_i\}\) 的递推公式

\[
\beta_1=c_1/b_1,\quad\beta_i=c_i/(b_i-a_i\beta_{i-1})\quad(i=2,3,\cdots,n-1);
\]

\textbf{步 2} 解 \(Ly=f\)：

\[
y_1=f_1/b_1,\quad y_i=(f_i-a_iy_{i-1})/(b_i-a_i\beta_{i-1})\quad(i=2,3,\cdots,n);
\]

\textbf{步 3} 解 \(Ux=y\)：

\[
x_n=y_n,\quad x_i=y_i-\beta_ix_{i+1}\quad(i=n-1,n-2,\cdots,2,1).
\]

\clearpage

原书 PDF 199；印刷页 186。

\href{/Users/xiangjunfeng/Desktop/数值计算专用/numerical-analysis/book/source-images/pdf-199.jpeg}{原页图像}

（续 7.4.3 追赶法。）

将计算系数 \(\beta_1\to\beta_2\to\cdots\to\beta_{n-1}\) 及
\(y_1\to y_2\to\cdots\to y_n\)
的过程称为\textbf{追的过程}，将计算方程组的解
\(x_n\to x_{n-1}\to\cdots\to x_1\) 的过程称为\textbf{赶的过程}。

总结上述讨论，有以下定理。

\textbf{定理 7.9} 设有三对角方程组 \(Ax=f\)，其中 \(A\)
满足对角占优条件，则 \(A\) 为非奇异矩阵且由追赶法计算公式中
\(\{\alpha_i\},\{\beta_i\}\) 满足：

\(1^\circ\) \(0<|\beta_i|<1\)（\(i=1,2,\cdots,n-1\)）；

\(2^\circ\)

\[
0<|c_i|\leqslant|b_i|-|a_i|<|\alpha_i|<|b_i|+|a_i|\quad(i=2,3,\cdots,n-1),
\]

\[
0<|b_n|-|a_n|<|\alpha_n|<|b_n|+|a_n|.
\]

追赶法公式实际上就是把 Gauss
消去法用到求解三对角方程组上去的结果。这时由于 \(A\)
特别简单，因此使得求解的计算公式也非常简单，计算量仅为 \(5n-4\)
次乘除法运算，而另外增加解一个方程组 \(Ax=f_2\) 仅需增加 \(3n-2\)
次乘除运算。易见追赶法的计算量是比较小的。

定理 7.9 的 \(1^\circ\)、\(2^\circ\)
说明追赶法计算公式中不会出现中间结果数量级的巨大增长和舍入误差的严重累积。

在用计算机计算时，只需用三个一维数组分别存储 \(A\) 的三条对角线元素
\(\{a_i\},\{b_i\},\{c_i\}\)，此外还需要用两组存储单元保存
\(\{\beta_i\},\{y_i\}\) 或 \(\{x_i\}\)。

\subsection{7.5
向量和矩阵的范数}\label{ux5411ux91cfux548cux77e9ux9635ux7684ux8303ux6570}

为了研究线性方程组近似解的误差估计和迭代法的收敛性，需要对
\(\mathbf R^n\)（\(n\) 维向量空间）中向量（或 \(\mathbf R^{n\times n}\)
中矩阵）的``大小''引进某种度量------向量（或矩阵）范数的概念。向量范数概念是三维
Euclid 空间中向量长度概念的推广，在数值分析中起着重要作用。

我们用 \(\mathbf R^n\) 表示 \(n\) 维实向量空间，用 \(\mathbf C^n\) 表示
\(n\) 维复向量空间，首先将向量长度概念推广到 \(\mathbf R^n\)（或
\(\mathbf C^n\)）中。

\textbf{定义 7.1} 设

\[
x=(x_1,x_2,\cdots,x_n)^{\mathrm T},\quad y=(y_1,y_2,\cdots,y_n)^{\mathrm T}\in\mathbf R^n\text{（或 }\mathbf C^n\text{）},
\]

将实数 \((x,y)=y^{\mathrm T}x=\sum_{i=1}^{n}x_iy_i\)（或复数
\((x,y)=y^Hx=\sum_{i=1}^{n}x_i\overline{y_i}\)，其中
\(y^H=\overline{y^{\mathrm T}}\)）称为向量 \(x,y\)
的\textbf{数量积}。将非负实数

\[
\|x\|_2=(x,x)^{\frac12}=\left(\sum_{i=1}^{n}x_i^2\right)^{\frac12}
\]

或

\[
\|x\|_2=(x,x)^{\frac12}=\left(\sum_{i=1}^{n}|x_i|^2\right)^{\frac12}
\]

（定义续下页。）

\clearpage

原书 PDF 200；印刷页 187。

\href{/Users/xiangjunfeng/Desktop/数值计算专用/numerical-analysis/book/source-images/pdf-200.jpeg}{原页图像}

（续 7.5 定义 7.1。）

称为向量 \(x\) 的 \textbf{Euclid 范数}。

下述定理可在线性代数书中找到。

\textbf{定理 7.10} 设 \(x,y\in\mathbf R^n\)（或 \(\mathbf C^n\)），则

\(1^\circ\) \((x,x)=0\)，当且仅当 \(x=0\) 时成立；

\(2^\circ\) \((\alpha x,y)=\alpha(x,y)\)，\(\alpha\) 为实数（或
\((x,\alpha y)=\overline\alpha(x,y)\)，\(\alpha\) 为复数）；

\(3^\circ\) \((x,y)=(y,x)\)（或 \((x,y)=\overline{(y,x)}\)）；

\(4^\circ\) \((x_1+x_2,y)=(x_1,y)+(x_2,y)\)；

\(5^\circ\)（Cauchy-Schwarz
不等式）\(|(x,y)|\leqslant\|x\|_2\cdot\|y\|_2\)，等式当且仅当 \(x\) 与
\(y\) 线性相关时成立；

\(6^\circ\)（三角不等式）\(\|x+y\|_2\leqslant\|x\|_2+\|y\|_2\)。

还可以用其他办法来度量 \(\mathbf R^n\) 中向量的``大小''。例如
\(x=(x_1,x_2)^{\mathrm T}\in\mathbf R^2\)，用一个 \(x\) 的函数
\(N(x)=\max_{i=1,2}|x_i|\) 求度量 \(x\) 的``大小''，而且这种度量
\(x\)``大小''的方法计算起来比 Euclid 范数方便。在许多应用中，对度量向量
\(x\)``大小''的函数 \(N(x)\)
都要求是正定的、齐次的且满足三角不等式。下面给出向量范数的一般定义。

\textbf{定义 7.2（向量的范数）} 如果向量 \(x\in\mathbf R^n\)（或
\(\mathbf C^n\)）的某个实值函数 \(N(x)=\|x\|\)，满足条件

\[
\begin{aligned}
1^\circ\quad &\|x\|\geqslant0\quad(\|x\|=0\text{ 当且仅当 }x=0)\quad\text{（正定条件）};\\
2^\circ\quad &\|\alpha x\|=|\alpha|\|x\|,\quad\forall\alpha\in\mathbf R\text{（或 }\alpha\in\mathbf C\text{）};\\
3^\circ\quad &\|x+y\|\leqslant\|x\|+\|y\|\quad\text{（三角不等式）},
\end{aligned}
\tag{7.5.1}
\]

则称 \(N(x)\) 是 \(\mathbf R^n\)（或
\(\mathbf C^n\)）上的一个\textbf{向量范数}（或\textbf{模}）。

由条件 \(3^\circ\) 可推出不等式

\[
\big|\|x\|-\|y\|\big|\leqslant\|x-y\|.
\tag{7.5.2}
\]

下面给出几种常用的向量范数。

（1）向量的
\textbf{\(\infty\)-范数（最大范数）}：\(\|x\|_\infty=\max_{1\leqslant i\leqslant n}|x_i|\)。

容易验证这样定义的向量 \(x\) 的函数 \(N(x)=\|x\|_\infty\)
满足向量范数的三个条件。

（2）向量的 \textbf{1-范数}：\(\|x\|_1=\sum_{i=1}^{n}|x_i|\)。同样可证
\(N(x)=\|x\|_1\) 是 \(\mathbf R^n\) 上的一个向量范数。

（3）向量的
\textbf{2-范数}：\(\|x\|_2=(x,x)^{\frac12}=\left(\sum_{i=1}^{n}x_i^2\right)^{\frac12}\)。由定理
7.10 知，\(N(x)=\|x\|_2\) 是 \(\mathbf R^n\) 上一个向量范数，称为向量
\(x\) 的 Euclid 范数。

（4）向量的
\textbf{\(p\)-范数}：\(\|x\|_p=\left(\sum_{i=1}^{n}|x_i|^p\right)^{1/p}\)，其中
\(p\in[1,\infty)\)。可以证明向量函数 \(N(x)=\|x\|_p\) 是 \(\mathbf R^n\)
上向量的范数，且容易说明上述三种范数是 \(p\)
范数的特殊情况（\(\|x\|_\infty=\lim_{p\to\infty}\|x\|_p\)）。

\textbf{例 7.6} 计算向量 \(x=(1,-2,3)^{\mathrm T}\) 的各种范数。

（例 7.6 解续下页。）

\clearpage

原书 PDF 201；印刷页 188。

\href{/Users/xiangjunfeng/Desktop/数值计算专用/numerical-analysis/book/source-images/pdf-201.jpeg}{原页图像}

（续 7.5 例 7.6。）

\textbf{解} \(\|x\|_1=6,\|x\|_\infty=3,\|x\|_2=\sqrt{14}\)。

\textbf{定义 7.3} 设 \(\{x^{(k)}\}\) 为 \(\mathbf R^n\)
中一向量序列，\(x^*\in\mathbf R^n\)，记
\(x^{(k)}=(x_1^{(k)},x_2^{(k)},\cdots,x_n^{(k)})^{\mathrm T}\)，\(x^*=(x_1^*,x_2^*,\cdots,x_n^*)^{\mathrm T}\)。如果
\(\lim_{k\to\infty}x_i^{(k)}=x_i^*\)（\(i=1,2,\cdots,n\)），则称
\(x^{(k)}\) \textbf{收敛}于向量 \(x^*\)，记作

\[
\lim_{k\to\infty}x^{(k)}=x^*.
\]

\textbf{定理 7.11（\(N(x)\) 的连续性）} 设非负函数 \(N(x)=\|x\|\) 为
\(\mathbf R^n\) 上任一向量范数，则 \(N(x)\) 是 \(x\) 分量
\(x_1,x_2,\cdots,x_n\) 的连续函数。

\textbf{证明} 设 \(x=\sum_{i=1}^{n}x_ie_i,y=\sum_{i=1}^{n}y_ie_i\)，其中
\(e_i=(0,\cdots,1,0,\cdots,0)^{\mathrm T}\)（即第 \(i\) 个元素为 1）。

只需证明当 \(x\to y\) 时 \(N(x)\to N(y)\) 即可。事实上，

\[
\begin{aligned}
|N(x)-N(y)|
&=\big|\|x\|-\|y\|\big|\leqslant\|x-y\|
=\left\|\sum_{i=1}^{n}(x_i-y_i)e_i\right\|\\
&\leqslant\sum_{i=1}^{n}|x_i-y_i|\|e_i\|
\leqslant\|x-y\|_\infty\sum_{i=1}^{n}\|e_i\|,
\end{aligned}
\]

即

\[
|N(x)-N(y)|\leqslant c\|x-y\|_\infty\to0\quad\text{（当 }x\to y\text{ 时）},
\]

其中

\[
c=\sum_{i=1}^{n}\|e_i\|.
\]

\textbf{定理 7.12（向量范数的等价性）} 设 \(\|x\|_s,\|x\|_t\) 为
\(\mathbf R^n\) 上向量的任意两种范数，则存在常数 \(c_1,c_2>0\)，使得

\[
c_1\|x\|_s\leqslant\|x\|_t\leqslant c_2\|x\|_s,\quad\text{对一切 }x\in\mathbf R^n.
\]

\textbf{证明} 只要就 \(\|x\|_s=\|x\|_\infty\)
证明上式成立即可，即证明存在常数 \(c_1,c_2>0\)，使

\[
c_1\leqslant\frac{\|x\|_t}{\|x\|_\infty}\leqslant c_2,\quad\text{对一切 }x\in\mathbf R^n\text{ 且 }x\neq0.
\]

考虑泛函 \(f(x)=\|x\|_t\geqslant0,\quad x\in\mathbf R^n\)。

记 \(S=\{x\mid\|x\|_\infty=1,x\in\mathbf R^n\}\)，则 \(S\)
是一个有界闭集。由于 \(f(x)\) 为 \(S\) 上的连续函数，所以 \(f(x)\) 于
\(S\) 上达到最大、最小值。设 \(x\in\mathbf R^n\) 且 \(x\neq0\)，则
\(\frac{x}{\|x\|_\infty}\in S\)，从而有

\[
f(x')=c_1\leqslant f\left(\frac{x}{\|x\|_\infty}\right)\leqslant c_2=f(x''),
\tag{7.5.3}
\]

其中 \(x',x''\in S\)。显然 \(c_1,c_2>0\)，上式为
\(c_1\leqslant\left\|\frac{x}{\|x\|_\infty}\right\|\leqslant c_2\)，即

\[
c_1\|x\|_\infty\leqslant\|x\|_t\leqslant c_2\|x\|_\infty,\quad\text{对一切 }x\in\mathbf R^n.
\]

注意，定理 7.12 不能推广到无穷维空间。由定理 7.12
可得到结论：如果在一种范数意义下向量序列收敛，则在任何一种范数意义下该向量序列亦收敛。

\clearpage

原书 PDF 202；印刷页 189。

\href{/Users/xiangjunfeng/Desktop/数值计算专用/numerical-analysis/book/source-images/pdf-202.jpeg}{原页图像}

（续 7.5 向量和矩阵的范数。）

\textbf{定理 7.13}
\(\lim_{k\to\infty}x^{(k)}=x^*\Longleftrightarrow\|x^{(k)}-x^*\|\to0\)（当
\(k\to\infty\) 时），其中 \(\|\cdot\|\) 为向量的任一种范数。

\textbf{证明}
显然，\(\lim_{k\to\infty}x^{(k)}=x^*\Longleftrightarrow\|x^{(k)}-x^*\|_\infty\to0\)（当
\(k\to\infty\) 时），而对于 \(\mathbf R^n\) 上任一种范数
\(\|\cdot\|\)，由定理 7.11，存在常数 \(c_1,c_2>0\)，使

\[
c_1\|x^{(k)}-x^*\|_\infty\leqslant\|x^{(k)}-x^*\|\leqslant c_2\|x^{(k)}-x^*\|_\infty,
\]

于是又有

\[
\|x^{(k)}-x^*\|_\infty\to0\Longleftrightarrow\|x^{(k)}-x^*\|\to0\quad\text{（当 }k\to\infty\text{ 时）}.
\]

下面将向量范数概念推广到矩阵上去。用 \(\mathbf R^{n\times n}\) 表示
\(n\times n\) 矩阵的集合（本质上是和 \(\mathbf R^{n\times n}\)
一样的向量空间），则由 \(\mathbf R^{n\times n}\) 上 2-范数可以得到
\(\mathbf R^{n\times n}\) 中矩阵的一种范数

\[
F(A)=\|A\|_F=\left(\sum_{i,j=1}^{n}a_{ij}^{2}\right)^{\frac12}
\]

称为 \(A\) 的 \textbf{Frobenius 范数}。\(\|A\|_F\)
显然满足正定性、齐次性及三角不等式。

下面给出矩阵范数的一般定义。

\textbf{定义 7.4（矩阵范数）} 如果矩阵 \(A\in\mathbf R^{n\times n}\)
的某个非负的实值函数 \(N(A)=\|A\|\)，满足条件

\[
\begin{aligned}
1^\circ\quad &\|A\|\geqslant0\quad(\|A\|=0\Longleftrightarrow A=0)\quad\text{（正定条件）};\\
2^\circ\quad &\|cA\|=|c|\|A\|,\quad c\text{ 为实数}\quad\text{（齐次条件）};\\
3^\circ\quad &\|A+B\|\leqslant\|A\|+\|B\|\quad\text{（三角不等式）};\\
4^\circ\quad &\|AB\|\leqslant\|A\|\|B\|,
\end{aligned}
\tag{7.5.4}
\]

则称 \(N(A)\) 是 \(\mathbf R^{n\times n}\)
上的一个\textbf{矩阵范数}（或\textbf{模}）。

上面定义的 \(F(A)=\|A\|_F\) 就是 \(\mathbf R^{n\times n}\)
上的一个矩阵范数。

由于在大多数与估计有关的问题中，矩阵和向量会同时参与讨论，所以希望引进一种矩阵的范数，它是和向量范数相联系而且和向量范数相容的，即

\[
\|Ax\|\leqslant\|A\|\|x\|
\tag{7.5.5}
\]

对于任意向量 \(x\in\mathbf R^n\) 及 \(A\in\mathbf R^{n\times n}\)
都成立。

为此再引进一种矩阵范数。

\textbf{定义 7.5（矩阵的算子范数）} 设
\(x\in\mathbf R^n,A\in\mathbf R^{n\times n}\)，给出一种向量范数
\(\|x\|_v\)（如 \(v=1,2\) 或 \(\infty\)），相应地定义一个矩阵的非负函数

\[
\|A\|_v=\max_{x\neq0}\frac{\|Ax\|_v}{\|x\|_v}.
\tag{7.5.6}
\]

可验证 \(\|A\|_v\) 满足定义 7.4（见下面定理），所以 \(\|A\|_v\) 是
\(\mathbf R^{n\times n}\) 上的一个矩阵范数，称为 \(A\)
的\textbf{算子范数}。

\textbf{定理 7.14} 设 \(\|x\|_v\) 是 \(\mathbf R^n\) 上一个向量范数，则
\(\|A\|_v\) 是 \(\mathbf R^{n\times n}\) 上的矩阵范数，且满足相容条件

\[
\|Ax\|_v\leqslant\|A\|_v\|x\|_v.
\tag{7.5.7}
\]

\textbf{证明}
由式（7.5.6）知，相容性条件式（7.5.7）是显然的。现只验证定义 7.4 中条

（证明接下页。）

\begin{quote}
校注（agent 补充，原书疑误）：定理 7.13 证明中原文引用``定理
7.11''，这里保留；所用的两侧常数比较直接对应上一页的定理
7.12（向量范数的等价性），定理 7.11
是连续性定理。\href{/Users/xiangjunfeng/Desktop/数值计算专用/numerical-analysis/book/staging/ch07b/assets/202-reference.png}{原文局部图}。
\end{quote}

\clearpage

原书 PDF 203；印刷页 190。

\href{/Users/xiangjunfeng/Desktop/数值计算专用/numerical-analysis/book/source-images/pdf-203.jpeg}{原页图像}

（续 7.5 定理 7.14 证明。）

件 \(4^\circ\)。

由式（7.5.7），有

\[
\|ABx\|_v\leqslant\|A\|_v\|Bx\|_v\leqslant\|A\|_v\|B\|_v\|x\|_v.
\]

当 \(x\neq0\) 时，有

\[
\frac{\|ABx\|_v}{\|x\|_v}\leqslant\|A\|_v\|B\|_v,\quad
\|AB\|_v=\max_{x\neq0}\frac{\|ABx\|_v}{\|x\|_v}\leqslant\|A\|_v\|B\|_v.
\]

显然这种矩阵范数 \(\|A\|_v\) 依赖于向量范数 \(\|x\|_v\)
的具体含义。也就是说，当给出一种具体的向量范数 \(\|x\|_v\)
时，相应地就得到了一种矩阵范数 \(\|A\|_v\)。

\textbf{定理 7.15} 设 \(x\in\mathbf R^n,A\in\mathbf R^{n\times n}\)，则

\(1^\circ\)
\(\|A\|_\infty=\max_{1\leqslant i\leqslant n}\sum_{j=1}^{n}|a_{ij}|\)（称为
\(A\) 的\textbf{行范数}）；

\(2^\circ\)
\(\|A\|_1=\max_{1\leqslant j\leqslant n}\sum_{i=1}^{n}|a_{ij}|\)（称为
\(A\) 的\textbf{列范数}）；

\(3^\circ\) \(\|A\|_2=\sqrt{\lambda_{\max}(A^{\mathrm T}A)}\)（称为
\(A\) 的 \textbf{2-范数}），

其中 \(\lambda_{\max}(A^{\mathrm T}A)\) 表示 \(A^{\mathrm T}A\)
的最大特征值。

\textbf{证明} 只就 \(1^\circ\)、\(3^\circ\) 给出证明，\(2^\circ\)
同理可证。

\(1^\circ\) 设 \(x=(x_1,x_2,\cdots,x_n)^{\mathrm T}\neq0\)，不妨设
\(A\neq0\)。记
\(t=\max_i|x_i|\)，\(\mu=\max_{1\leqslant i\leqslant n}\sum_{j=1}^{n}|a_{ij}|\)，则

\[
\|Ax\|_\infty
=\max_{1\leqslant i\leqslant n}\left|\sum_{j=1}^{n}a_{ij}x_j\right|
\leqslant\max_i\sum_{j=1}^{n}|a_{ij}||x_j|
\leqslant t\max_i\sum_{j=1}^{n}|a_{ij}|.
\]

这说明，对于任何非零 \(x\in\mathbf R^n\)，有

\[
\frac{\|Ax\|_\infty}{\|x\|_\infty}\leqslant\mu.
\tag{7.5.8}
\]

下面说明，有一向量 \(x_0\neq0\)，使
\(\frac{\|Ax_0\|_\infty}{\|x_0\|_\infty}=\mu\)。设
\(\mu=\sum_{j=1}^{n}|a_{i_0j}|\)，取向量

\[
x_0=(x_1,x_2,\cdots,x_n)^{\mathrm T},
\]

其中

\[
x_j=\operatorname{sgn}(a_{i_0j})\quad(j=1,2,\cdots,n).
\]

显然 \(\|x_0\|_\infty=1\)，且 \(Ax_0\) 的第 \(i_0\) 个分量为
\(\sum_{i=1}^{n}a_{i_0j}x_j=\sum_{j=1}^{n}|a_{i_0j}|\)，这说明

\[
\|Ax_0\|_\infty=\max_i\sum_{j=1}^{n}|a_{ij}x_j|=\sum_{j=1}^{n}|a_{i_0j}|=\mu.
\]

\(3^\circ\) 由于
\(\|Ax\|_2^2=(Ax,Ax)=(A^{\mathrm T}Ax,x)\geqslant0\)，对于一切
\(x\in\mathbf R^n\)，从而 \(A^{\mathrm T}A\) 的特征值为非负实数，设为

\[
\lambda_1\geqslant\lambda_2\geqslant\cdots\geqslant\lambda_n\geqslant0.
\tag{7.5.9}
\]

\(A^{\mathrm T}A\) 为对称矩阵，设 \(u_1,u_2,\cdots,u_n\) 为
\(A^{\mathrm T}A\) 的相应于式（7.5.9）的特征向量且

（证明续下页。）

\begin{quote}
校注（agent 补充，原书疑误）：\(Ax_0\) 第 \(i_0\)
个分量表达式的第一个求和号下限原图为 \(i=1\)，这里保留；该和式的项是
\(a_{i_0j}x_j\)，求和指标应为 \(j\)，相应下限应为
\(j=1\)。\href{/Users/xiangjunfeng/Desktop/数值计算专用/numerical-analysis/book/staging/ch07b/assets/203-sum-index.png}{原式局部图}。
\end{quote}

\clearpage

原书 PDF 204；印刷页 191。

\href{/Users/xiangjunfeng/Desktop/数值计算专用/numerical-analysis/book/source-images/pdf-204.jpeg}{原页图像}

（续 7.5 定理 7.15 证明。）

\((u_i,u_j)=\delta_{ij}\)，又设 \(x\in\mathbf R^n\)
为任一非零向量，于是有

\[
x=\sum_{i=1}^{n}c_iu_i,
\]

其中 \(c_i\) 为组合系数。

\[
\frac{\|Ax\|_2^2}{\|x\|_2^2}
=\frac{(A^{\mathrm T}Ax,x)}{(x,x)}
\leqslant\frac{\sum_{i=1}^{n}c_i^2\lambda_1}{\sum_{i=1}^{n}c_i^2}=\lambda_1.
\]

另一方面，取 \(x=u_1\)，则上式等号成立，故

\[
\|A\|_2=\max_{x\neq0}\frac{\|Ax\|_2}{\|x\|_2}=\sqrt{\lambda_1}=\sqrt{\lambda_{\max}(A^{\mathrm T}A)}.
\]

由定理 7.15 看出，计算一个矩阵的 \(\|A\|_\infty,\|A\|_1\)
还是比较容易的，而矩阵的 2-范数 \(\|A\|_2\) 在计算上不方便。但是矩阵的
2-范数具有许多好的性质，它在理论上是有用的。

\textbf{例 7.7} \(A=\begin{pmatrix}1&-2\\-3&4\end{pmatrix}\)，计算 \(A\)
的各种范数。

\textbf{解}
\(\|A\|_1=6,\|A\|_\infty=7,\|A\|_F\approx5.477,\|A\|_2=\sqrt{15+\sqrt{221}}\approx5.46\)。

对于复矩阵（即 \(A\in\mathbf C^{n\times n}\)），定理 7.15 中的
\(1^\circ\)、\(2^\circ\) 显然成立；\(3^\circ\) 则应改为

\[
\|A\|_2=\max_{x\neq0}\left(\frac{x^HA^HAx}{x^Hx}\right)=\sqrt{\lambda_{\max}(A^HA)}.
\]

其中 \(A^H=\overline A^{\mathrm T}\)（\(A\) 共轭转置）。

\textbf{定义 7.6} 设 \(A\in\mathbf R^{n\times n}\) 的特征值为
\(\lambda_i\)（\(i=1,2,\cdots,n\)），称
\(\rho(A)=\max_{1\leqslant i\leqslant n}|\lambda_i|\) 为 \(A\)
的\textbf{谱半径}。

\textbf{定理 7.16（特征值上界）} 设 \(A\in\mathbf R^{n\times n}\)，则
\(\rho(A)\leqslant\|A\|\)，即 \(A\) 的谱半径不超过 \(A\)
的任何一种算子范数（对于 \(\|A\|_F\) 亦对）。

\textbf{证明} 设 \(\lambda\) 是 \(A\) 的任一特征值，\(x\)
为相应的特征向量，则 \(Ax=\lambda x\)，由式（7.5.7）得

\[
|\lambda|\|x\|=\|\lambda x\|=\|Ax\|\leqslant\|A\|\|x\|,
\]

即

\[
|\lambda|\leqslant\|A\|.
\]

\textbf{定理 7.17} 如果 \(A\in\mathbf R^{n\times n}\) 为对称矩阵，则
\(\|A\|_2=\rho(A)\)。

证明留作习题。

\textbf{定理 7.18} 如果 \(\|B\|<1\)，则 \(I\pm B\) 为非奇异矩阵，且
\(\|(I\pm B)^{-1}\|\leqslant\frac1{1-\|B\|}\)，其中 \(\|\cdot\|\)
是指矩阵的算子范数。

\textbf{证明} 用反证法。若 \(\det(I\pm B)=0\)，则 \((I\pm B)x=0\)
有非零解，即存在 \(x_0\neq0\) 使

（证明续下页。）

\begin{quote}
校注（agent 补充，原书疑误）：复矩阵 2-范数公式的 Rayleigh
商外，原图没有平方根或 \(1/2\) 次幂，此处照录。该最大商等于
\(\lambda_{\max}(A^HA)=\|A\|_2^2\)；要与两侧 \(\|A\|_2\) 和
\(\sqrt{\lambda_{\max}(A^HA)}\)
一致，中间项需开平方。\href{/Users/xiangjunfeng/Desktop/数值计算专用/numerical-analysis/book/staging/ch07b/assets/204-complex-norm.png}{原式局部图}。
\end{quote}

\clearpage

原书 PDF 205；印刷页 192。

\href{/Users/xiangjunfeng/Desktop/数值计算专用/numerical-analysis/book/source-images/pdf-205.jpeg}{原页图像}

（续 7.5 定理 7.18 证明。）

\(Bx_0=x_0\)，\(\frac{\|Bx_0\|}{\|x_0\|}=1\)，故
\(\|B\|\geqslant1\)，与假设矛盾。又由 \((I\pm B)(I\pm B)^{-1}=I\)，有

\[
(I\pm B)^{-1}=I\mp B(I-B)^{-1},
\]

从而

\[
\|(I\pm B)^{-1}\|\leqslant\|I\|+\|B\|\|(I\pm B)^{-1}\|,
\]

\[
\|(I\pm B)^{-1}\|\leqslant\frac1{1-\|B\|}.
\]

\subsection{7.6 误差分析}\label{ux8befux5deeux5206ux6790}

\subsubsection{7.6.1
矩阵的条件数}\label{ux77e9ux9635ux7684ux6761ux4ef6ux6570}

考虑线性方程组 \(Ax=b\)，其中设 \(A\) 为非奇异矩阵，\(x\)
为方程组的精确解。

由于 \(A\)（或 \(b\)）元素是测量得到的，或者是计算的结果，所以 \(A\)（或
\(b\)）常带有某些观测误差或者包含有舍入误差。因此，在处理实际问题
\((A+\delta A)x=b+\delta b\) 时，通常要考虑 \(A\)（或
\(b\)）的微小误差对解的影响，即考虑估计 \(x-y\)，其中 \(y\) 是
\((A+\delta A)y=b+\delta b\) 的解。

首先考察一个例子。

\textbf{例 7.8} 设有方程组

\[
\begin{pmatrix}1&1\\1&1.0001\end{pmatrix}
\begin{pmatrix}x_1\\x_2\end{pmatrix}
=\begin{pmatrix}2\\2\end{pmatrix}
\]

或

\[
Ax=b,
\tag{7.6.1}
\]

它的精确解为

\[
x=(2,0)^{\mathrm T}.
\]

现在考虑常数项的微小变化对方程组解的影响，即考察方程组

\[
\begin{pmatrix}1&1\\1&1.000\,1\end{pmatrix}
\begin{pmatrix}y_1\\y_2\end{pmatrix}
=\begin{pmatrix}2\\2.000\,1\end{pmatrix}
\]

或

\[
A(x+\delta x)=b+\delta b,
\tag{7.6.2}
\]

其中 \(\delta b=(0,0.000\,1)^{\mathrm T},y=x+\delta x\)，\(x\)
为式（7.6.1）的解。显然方程组（7.6.2）的解为

\[
x+\delta x=(1,1)^{\mathrm T}.
\]

可以看到，式（7.6.1）的常数项 \(b\) 的第二个分量只有 \(\frac1{10\,000}\)
的微小变化，方程组的解却变化很大。这样的方程组称为病态方程组。

\textbf{定义 7.7} 如果矩阵 \(A\) 或常数项 \(b\) 的微小变化，引起方程组
\(Ax=b\) 解的巨大变化，则称此方程组为\textbf{病态方程组}，矩阵 \(A\)
称为\textbf{病态矩阵}（相对于方程组而言），否则称方程组为\textbf{良态方程组}，\(A\)
称为\textbf{良态矩阵}。

应该注意，矩阵的病态性质是矩阵本身的特性，下面希望找出刻画矩阵病态性质

（正文续下页。）

\begin{quote}
校注（agent 补充，原书疑误）：本页定理 7.18 证明原文为
\(Bx_0=x_0\)，以及
\((I\pm B)^{-1}=I\mp B(I-B)^{-1}\)，均按印刷内容保留。若同时处理 \(I+B\)
和 \(I-B\) 两种情形，前式应为 \(Bx_0=\mp x_0\)，后式右端的逆矩阵应为
\((I\pm B)^{-1}\)；原文写法只与减号分支一致。前式符号不影响随后所取范数等于
1
的结论。\href{/Users/xiangjunfeng/Desktop/数值计算专用/numerical-analysis/book/staging/ch07b/assets/205-inverse-signs.png}{原式局部图}。
\end{quote}

\clearpage

原书 PDF 206；印刷页 193。

\href{/Users/xiangjunfeng/Desktop/数值计算专用/numerical-analysis/book/source-images/pdf-206.jpeg}{原页图像}

（续 7.6.1 矩阵的条件数。）

的量。设有方程组

\[
Ax=b,
\tag{7.6.3}
\]

其中 \(A\) 为非奇异阵，\(x\)
为式（7.6.3）的精确解。下面研究当方程组的系数矩阵 \(A\)（或
\(b\)）有微小误差（扰动）时对解的影响。

现设 \(A\) 是精确的，\(b\) 有误差 \(\delta b\)，解为 \(x+\delta x\)，则

\[
A(x+\delta x)=b+\delta b,\quad\delta x=A^{-1}\delta b,\quad\|\delta x\|\leqslant\|A^{-1}\|\|\delta b\|.
\tag{7.6.4}
\]

由式（7.6.3），有

\[
\|b\|\leqslant\|A\|\|x\|,
\]

即

\[
\frac1{\|x\|}\leqslant\frac{\|A\|}{\|b\|}\quad\text{（设 }b\neq0\text{）},
\tag{7.6.5}
\]

于是由式（7.6.4）式及式（7.6.5），有如下定理。

\textbf{定理 7.19} 设 \(A\) 是非奇异阵，\(Ax=b\neq0\)，且
\(A(x+\delta x)=b+\delta b\)，则

\[
\frac{\|\delta x\|}{\|x\|}\leqslant\|A^{-1}\|\|A\|\frac{\|\delta b\|}{\|b\|}.
\]

上式给出了解的相对误差的上界，常数项 \(b\) 的相对误差在解中可能放大
\(\|A^{-1}\|\|A\|\) 倍。

现设 \(b\) 是精确的，\(A\) 有微小误差（扰动）\(\delta A\)，解为
\(x+\delta x\)，则

\[
(A+\delta A)(x+\delta x)=b,
\]

\[
(A+\delta A)\delta x=-(\delta A)x.
\tag{7.6.6}
\]

如果 \(\delta A\) 不受限制的话，\(A+\delta A\) 可能奇异，而

\[
(A+\delta A)=A(I+A^{-1}\delta A).
\]

由定理 7.18 知，当 \(\|A^{-1}\delta A\|<1\)
时，\((I+A^{-1}\delta A)^{-1}\) 存在，由式（7.6.6），有

\[
\delta x=-(I+A^{-1}\delta A)^{-1}A^{-1}(\delta A)x,\quad
\|\delta x\|\leqslant\frac{\|A^{-1}\|\|\delta A\|\|x\|}{1-\|A^{-1}(\delta A)\|}.
\]

设 \(\|A^{-1}\|\|\delta A\|<1\)，即得

\[
\frac{\|\delta x\|}{\|x\|}
\leqslant
\frac{\|A^{-1}\|\|A\|\dfrac{\|\delta A\|}{\|A\|}}
{1-\|A^{-1}\|\|A\|\dfrac{\|\delta A\|}{\|A\|}}.
\tag{7.6.7}
\]

\textbf{定理 7.20} 设 \(A\) 为非奇异矩阵，\(Ax=b\neq0\)，且
\((A+\delta A)(x+\delta x)=b\)，如果
\(\|A^{-1}\|\|\delta A\|<1\)，则式（7.6.7）成立。

如果 \(\delta A\) 充分小，且在 \(\|A^{-1}\|\|\delta A\|<1\)
条件下，那么式（7.6.7）说明矩阵 \(A\) 的相对误差
\(\frac{\|\delta A\|}{\|A\|}\) 在解中可能放大 \(\|A^{-1}\|\|A\|\) 倍。

总之，量 \(\|A^{-1}\|\|A\|\) 愈小，由 \(A\)（或
\(b\)）的相对误差引起的解的相对误差就愈小；量 \(\|A^{-1}\|\|A\|\)
愈大，该解的相对误差就可能愈大。所以量 \(\|A^{-1}\|\|A\|\)
实际上刻画了解对原始数据变化的灵敏度，即刻画了方程组的病态程度，于是引进下述定

（定义续下页。）

\clearpage

原书 PDF 207；印刷页 194。

\href{/Users/xiangjunfeng/Desktop/数值计算专用/numerical-analysis/book/source-images/pdf-207.jpeg}{原页图像}

（续 7.6.1 矩阵的条件数。）

义。

\textbf{定义 7.8} 设 \(A\) 为非奇异矩阵，称数
\(\operatorname{cond}(A)_v=\|A^{-1}\|_v\|A\|_v\)（\(v=1,2\) 或
\(\infty\)）为矩阵 \(A\) 的\textbf{条件数}。

由此看出矩阵的条件数与范数有关。

矩阵的条件数是一个十分重要的概念。由上面讨论知，当 \(A\)
的条件数相对地大时，即 \(\operatorname{cond}(A)\gg1\)
时，式（7.6.3）是病态的（即 \(A\) 是病态矩阵，或者说 \(A\)
是坏条件的）；当 \(A\) 的条件数相对地小时，式（7.6.3）是良态的（或者说
\(A\) 是好条件的）。\(A\)
的条件数愈大，方程组的病态程度愈严重，也就愈难得到方程组的比较准确的解。

通常使用的条件数有

（1）\(\operatorname{cond}(A)_\infty=\|A^{-1}\|_\infty\|A\|_\infty\)；

（2）\(A\) 的\textbf{谱条件数}

\[
\operatorname{cond}(A)_2=\|A\|_2\|A^{-1}\|_2
=\sqrt{\frac{\lambda_{\max}(A^{\mathrm T}A)}{\lambda_{\min}(A^{\mathrm T}A)}}.
\]

当 \(A\)
为对称阵时，\(\operatorname{cond}(A)_2=\frac{|\lambda_1|}{|\lambda_n|}\)，其中
\(\lambda_1,\lambda_n\) 为 \(A\) 的绝对值最大和绝对值最小的特征值。

条件数有如下性质：

（1）对任何非奇异矩阵 \(A\)，都有
\(\operatorname{cond}(A)_v\geqslant1\)。事实上，

\[
\operatorname{cond}(A)_v=\|A^{-1}\|_v\|A\|_v\geqslant\|A^{-1}A\|_v=1.
\]

（2）设 \(A\) 为非奇异矩阵，\(c\) 为不等于零的常数，则
\(\operatorname{cond}(cA)_v=\operatorname{cond}(A)_v\)。

（3）如果 \(A\) 为正交矩阵，则 \(\operatorname{cond}(A)_2=1\)；如果
\(A\) 为非奇异矩阵，\(R\) 为正交矩阵，则

\[
\operatorname{cond}(RA)_2=\operatorname{cond}(AR)_2=\operatorname{cond}(A)_2.
\]

\textbf{例 7.9} 已知 Hilbert 矩阵

\[
H_n=\begin{pmatrix}
1&\frac12&\cdots&\frac1n\\[4pt]
\frac12&\frac13&\cdots&\frac1{n+1}\\
\vdots&\vdots&&\vdots\\
\frac1n&\frac1{1+n}&\cdots&\frac1{2n-1}
\end{pmatrix},
\]

计算 \(H_3\) 的条件数。

\textbf{解}

\[
H_3=\begin{pmatrix}
1&\frac12&\frac13\\[4pt]
\frac12&\frac13&\frac14\\[4pt]
\frac13&\frac14&\frac15
\end{pmatrix},\quad
H_3^{-1}=\begin{pmatrix}
9&-36&30\\
-36&192&-180\\
30&-180&180
\end{pmatrix}.
\]

（例 7.9 解续下页。）

\clearpage

原书 PDF 208；印刷页 195。

\href{/Users/xiangjunfeng/Desktop/数值计算专用/numerical-analysis/book/source-images/pdf-208.jpeg}{原页图像}

（续 7.6.1 例 7.9。）

（1）计算 \(H_3\) 条件数 \(\operatorname{cond}(H_3)_\infty\)。

\[
\|H_3\|_\infty=11/6,\quad\|H_3^{-1}\|_\infty=408,
\]

所以 \(\operatorname{cond}(H_3)_\infty=748\)。同样可计算
\(\operatorname{cond}(H_6)=2.9\times10^6\)。一般当 \(n\) 愈大时，\(H_n\)
的病态愈严重。

（2）考虑方程组

\[
H_3x=(11/6,13/12,47/60)^{\mathrm T}=b.
\]

设 \(H_3\) 及 \(b\) 有微小误差（取三位有效数字），有

\[
\begin{pmatrix}
1.00&0.500&0.333\\
0.500&0.333&0.250\\
0.333&0.250&0.200
\end{pmatrix}
\begin{pmatrix}x_1+\delta x_1\\x_2+\delta x_2\\x_3+\delta x_3\end{pmatrix}
=\begin{pmatrix}1.83\\1.08\\0.783\end{pmatrix},
\tag{7.6.8}
\]

简记作 \((H_3+\delta H_3)(x+\delta x)=b+\delta b\)。方程组 \(H_3x=b\)
与式（7.6.8）的精确解为

\[
x=(1,1,1)^{\mathrm T},\quad
x+\delta x=(1.089\,512\,538,\ 0.487\,967\,062,\ 1.491\,002\,798)^{\mathrm T},
\]

于是

\[
\delta x=(0.089\,5,\ -0.512\,0,\ 0.491\,0)^{\mathrm T},\quad
\frac{\|\delta H_3\|_\infty}{\|H_3\|_\infty}\approx0.18\times10^{-3}<0.02\%,
\]

\[
\frac{\|\delta b\|_\infty}{\|b\|_\infty}\approx0.182\%,\quad
\frac{\|\delta x\|_\infty}{\|x\|_\infty}\approx51.2\%.
\]

这就是说，\(H_3\) 与 \(b\) 相对误差不超过
\(0.3\%\)，而引起解的相对误差超过 \(50\%\)。

由上面讨论可知，要判别一个矩阵是否病态，需要计算条件数
\(\operatorname{cond}(A)=\|A^{-1}\|\|A\|\)，而计算 \(A^{-1}\)
是比较费劲的。那么在实际计算中如何发现病态情况呢？

（1）如果在 \(A\)
的三角约化时（尤其是用主元素消去法解方程组（7.6.3）时）出现小主元，那么对大多数矩阵来说，\(A\)
是病态矩阵。例如用选主元的直接三角分解法解方程组（7.6.8）（用双精度累加计算
\(\sum_i a_ib_i\)，结果舍入为三位浮点数），则有

\[
I_{23}(H_3+\delta H_3)
=\begin{pmatrix}
1&&\\
0.333&1&\\
0.500&0.994&1
\end{pmatrix}
\begin{pmatrix}
1&0.500\,0&0.333\,0\\
&0.083\,5&0.089\,1\\
&&\boxed{-0.005\,07}
\end{pmatrix}
=LU.
\]

（2）如果 \(A\) 的最大特征值和最小特征值之比（按绝对值）是大的，则 \(A\)
是病态的，即当 \(|\lambda_1|/|\lambda_n|\) 是大的，其中 \(A\)
特征值次序为

\[
|\lambda_1|\geqslant|\lambda_2|\geqslant\cdots\geqslant|\lambda_n|>0.
\]

显然

\[
|\lambda_1|\leqslant\|A\|,\quad\frac1{|\lambda_n|}\leqslant\|A^{-1}\|,
\]

因此

\[
\operatorname{cond}(A)\geqslant\frac{|\lambda_1|}{|\lambda_n|}\gg1.
\]

（3）如果系数矩阵的行列式值相对来说很小，或系数矩阵某些行近似线性相关，则
\(A\) 可能是病态的。

（条目（4）续下页。）

\begin{quote}
校注（agent 补充，原书疑误）：原文印作
\(\operatorname{cond}(H_6)=2.9\times10^6\)，这里保留。沿用本例的
\(\infty\)-范数，对 \(6\times6\) Hilbert 矩阵作有理数精确运算得到
\(\|H_6\|_\infty\|H_6^{-1}\|_\infty=29\,070\,279\approx2.9\times10^7\)，原书指数疑应为
7。\href{/Users/xiangjunfeng/Desktop/数值计算专用/numerical-analysis/book/staging/ch07b/assets/208-h6-cond.png}{原文局部图}。
\end{quote}

\begin{quote}
校注（agent
补充，ch07b-E014）：本页称为``精确解''的小数向量按原书照录。对式（7.6.8）所列有限小数作精确有理数求解，结果为
\((153\,971,68\,960,210\,710)^{\mathrm T}/141\,321\approx(1.0895125282,0.4879671103,1.4910027526)^{\mathrm T}\)，与所印三个分量的末几位不符。因此原书这些小数不能视作式（7.6.8）的精确解；该差异不影响后文约
\(51.2\%\) 的结论。
\end{quote}

\clearpage

原书 PDF 209；印刷页 196。

\href{/Users/xiangjunfeng/Desktop/数值计算专用/numerical-analysis/book/source-images/pdf-209.jpeg}{原页图像}

（续 7.6.1 矩阵的条件数。）

（4）如果系数矩阵 \(A\) 元素间数量级相差很大，并且无一定规则，则 \(A\)
可能是病态的。

病态问题通常不能用选主元素的消去法来解决。对于此类问题一般采用高精度的算术运算（采用双倍字长进行运算）或者预处理方法。即将求解
\(Ax=b\) 的问题转化为求解一等价方程组

\[
\begin{cases}
PAQy=Pb,\\
y=Q^{-1}x.
\end{cases}
\]

通过选择非奇异矩阵 \(P,Q\)，使
\(\operatorname{cond}(PAQ)<\operatorname{cond}(A)\)。一般选择 \(P,Q\)
为对角阵或者三角阵。

当矩阵 \(A\) 的元素大小不均时，在 \(A\)
的行（或列）中引进适当的比例因子（使矩阵 \(A\) 的所有行或列按
\(\infty\)-范数大体上有相同的长度，使 \(A\) 的系数均衡），对 \(A\)
的条件数是有影响的。但这种方法不能保证 \(A\) 的条件数一定得到改善。

\textbf{例 7.10} 设 \(A=\begin{pmatrix}1&10^4\\1&1\end{pmatrix}\)，且

\[
\begin{pmatrix}1&10^4\\1&1\end{pmatrix}
\begin{pmatrix}x_1\\x_2\end{pmatrix}
=\begin{pmatrix}10^4\\2\end{pmatrix},
\tag{7.6.9}
\]

计算 \(\operatorname{cond}(A)_\infty\)。

\textbf{解} 由

\[
A=\begin{pmatrix}1&10^4\\1&1\end{pmatrix},
\]

有

\[
A^{-1}=\frac1{10^4-1}\begin{pmatrix}-1&10^4\\1&-1\end{pmatrix},\quad
\operatorname{cond}(A)_\infty=\frac{(1+10^4)^2}{10^4-1}\approx10^4.
\]

现在 \(A\) 的第 1 行引进比例因子。如用
\(s_1=\max_{1\leqslant i\leqslant2}|a_{1i}|=10^4\) 除第一个方程式，得
\(A'x=b'\)，即

\[
\begin{pmatrix}10^{-4}&1\\1&1\end{pmatrix}
\begin{pmatrix}x_1\\x_2\end{pmatrix}
=\begin{pmatrix}1\\2\end{pmatrix},
\tag{7.6.10}
\]

\[
(A')^{-1}=\frac1{1-10^{-4}}\begin{pmatrix}-1&1\\1&-10^{-4}\end{pmatrix},\quad
\operatorname{cond}(A')_\infty=\frac1{1-10^{-4}}\approx4.
\]

当用列主元素消去法解式（7.6.9）时（计算到三位数字），得

\[
(A,b)\longrightarrow
\left(\begin{array}{cc|c}
1&10^4&10^4\\
0&-10^4&-10^4
\end{array}\right),
\]

于是得到很坏的结果：\(x_1=0,x_2=1\)。

现用列主元素消去法解式（7.6.10），得

（例 7.10 解续下页。）

\begin{quote}
校注（agent 补充，原书疑误）：\(\operatorname{cond}(A')_\infty\)
的分式分子原图为 1，而后面印作
\(\approx4\)，这里保留。由本页两个矩阵，\(\|A'\|_\infty=2\)，\(\|(A')^{-1}\|_\infty=2/(1-10^{-4})\)，故条件数应为
\(4/(1-10^{-4})=40\,000/9\,999\approx4.0004\)。\href{/Users/xiangjunfeng/Desktop/数值计算专用/numerical-analysis/book/staging/ch07b/assets/209-scaled-cond.png}{原式局部图}。
\end{quote}

\clearpage

原书 PDF 210；印刷页 197。

\href{/Users/xiangjunfeng/Desktop/数值计算专用/numerical-analysis/book/source-images/pdf-210.jpeg}{原页图像}

（续 7.6.1 例 7.10。）

\[
(A',b')\longrightarrow
\left(\begin{array}{cc|c}1&1&2\\10^{-4}&1&1\end{array}\right)
\longrightarrow
\left(\begin{array}{cc|c}1&1&2\\0&1&1\end{array}\right),
\]

从而得到较好的结果：\(x_1=1,x_2=1\)。

设 \(\overline x\) 为方程组 \(Ax=b\) 的近似解，于是可计算
\(\overline x\) 的剩余向量 \(r=b-A\overline x\)，当 \(r\)
很小时，\(\overline x\) 是否为 \(Ax=b\)
一个较好的近似解呢？下述定理给出了解答。

\textbf{定理 7.21（事后误差估计）}

\(1^\circ\) 设 \(A\) 为非奇异矩阵，\(x\) 是精确解，\(Ax=b\neq0\)；

\(2^\circ\) 设 \(\overline x\)
是方程组的近似解，\(r=b-A\overline x\)，则

\[
\frac{\|x-\overline x\|}{\|x\|}\leqslant\operatorname{cond}(A)\cdot\frac{\|r\|}{\|b\|}.
\tag{7.6.11}
\]

\textbf{证明} 由 \(x-\overline x=A^{-1}r\)，得

\[
\|x-\overline x\|\leqslant\|A^{-1}\|\|r\|,
\tag{7.6.12}
\]

又有

\[
\|b\|=\|Ax\|\leqslant\|A\|\|x\|,
\]

\[
\frac1{\|x\|}\leqslant\frac{\|A\|}{\|b\|},
\tag{7.6.13}
\]

由式（7.6.12）及式（7.6.13）即得到式（7.6.11）。

\subsubsection{7.6.2 舍入误差}\label{ux820dux5165ux8befux5dee}

在复杂的计算中，由浮点运算而引进的舍入误差可能积累而影响答案，因此对任何算法都需要进行舍入误差分析，看其是否过度影响所得的结果。

下面叙述采用选主元素 Gauss
消去法解式（7.6.3）的计算过程中舍入误差对解的影响。在 Gauss
消去法的舍入误差的分析和研究方面，von Neumann（在 1947 年）、Givens（在
1954 年）和 Wilkinson（在 1961 年、1963 年）等人都作出了自己的贡献。

设 \(\overline x\) 为用选主元素 Gauss 消去法解式（7.6.3）的计算解，\(x\)
为式（7.6.3）的精确解。若要直接计算每一步舍入误差对解的影响来获得界的估计
\(\|x-\overline x\|\)，那将是非常困难的。Wilkinson
等人提出了``向后误差分析方法''，其基本思想是把计算过程中舍入误差对解的影响归结为原始数据变化对解的影响，即计算解
\(\overline x\) 是下述扰动方程组的精确解 \((A+\delta A)x=b\)，其中
\(\delta A\) 为某个``小''矩阵。

下面给出一个定理来说明这个结果。

\textbf{定理 7.22（选主元素 Gauss 消去法误差分析）} 如果

\(1^\circ\) 设 \(A\) 为 \(n\) 阶非奇异矩阵；

\(2^\circ\) 用列主元素消去法（或完全主元素消去法）解式（7.6.3）；

（定理条件与结论续下页。）

\begin{quote}
校注（agent 补充，原书疑误）：向后误差分析段中印刷方程为
\((A+\delta A)x=b\)，而紧邻文字明确指计算解 \(\overline x\)
是其精确解。此处保留原式；按本段符号约定，应写
\((A+\delta A)\overline x=b\)。\href{/Users/xiangjunfeng/Desktop/数值计算专用/numerical-analysis/book/staging/ch07b/assets/210-backward-error.png}{原文局部图}。
\end{quote}

\clearpage

原书 PDF 211；印刷页 198。

\href{/Users/xiangjunfeng/Desktop/数值计算专用/numerical-analysis/book/source-images/pdf-211.jpeg}{原页图像}

（续 7.6.2 定理 7.22。）

\(3^\circ\) 记
\(a_k=\max_{1\leqslant i,j\leqslant n}|a_{ij}^{(k)}|\)，\(a=\max_{1\leqslant i,j\leqslant n}|a_{ij}|\)，\(r=\max_{1\leqslant k\leqslant n}a_k/a\)------元素的增长因子，\(A^{(k)}\equiv(a_{ij}^{(k)})_n\)；

\(4^\circ\) \(t\) 为计算机字长（指尾数部分），矩阵的阶数 \(n\) 满足
\(n\cdot2^{-t}\leqslant0.01\)，

则（1）用选主元素 Gauss 消去法计算的三角阵 \(L,U\) 满足 \(LU=A+E\)，其中

\[
|(E)_{ij}|\leqslant2(n-1)ra2^{-t};
\]

（2）用选主元素 Gauss 消去法得到的计算解 \(\overline x\) 精确满足
\((A+\delta A)x=b\)，其中

\[
\|\delta A\|_\infty\leqslant1.01(n^3+3n^2)ar2^{-t}\leqslant1.01(n^3+3n^2)r\|A\|_\infty2^{-t};
\]

（3）计算解精度估计（\(x\) 为 \(Ax=b\) 精确解）

\[
\frac{\|x-\overline x\|_\infty}{\|x\|_\infty}
\quad
\frac{\operatorname{cond}(A)_\infty}
{1-\operatorname{cond}(A)_\infty\dfrac{\|\delta A\|_\infty}{\|A\|_\infty}}
\times1.01(n^3+3n^2)r2^{-t}.
\tag{7.6.14}
\]

上述计算解 \(\overline x\) 精度估计式说明，\(\overline x\)
的相对误差限依赖于
\(\operatorname{cond}(A)_\infty\)、元素的增长因子、方程组阶数、计算机字长等。

\subsection{小结}\label{ux5c0fux7ed3}

在 Gauss
消去法中引进选主元素的技巧，就得到了解方程组的完全主元素消去法和列主元素消去法，对一般矩阵引进选主元素的技巧的根本作用是为了对增长因子进行控制。完全选主元素及列主元素方法都是数值稳定的算法（即对舍入误差的增长得以控制的算法）。用完全主元素消去法解方程组（至少对于良态方程组是这样）具有较高的精确度，但它需要花费较多的机器时间。列主元素消去法是比完全主元素消去法更实用的算法，一般使用较多。这两种方法都是计算机上解线性方程组的有效方法。用
Gauss-Jordan 方法求逆矩阵是比较方便的。

从代数角度看，直接分解法和 Gauss
消去法本质上一样，但如果采用``双精度累加''计算
\(\sum_i a_ib_i\)，那么直接三角分解法的精度要比 Gauss
消去法的精度高，如果不采用``双精度累加''（\(\sum_i a_ib_i\)
中每一个乘积分别进行舍入），那么两个方法将给出相同的结果。

对于对称正定矩阵
\(A\)，采用不选主元素的平方根法（或改进的平方根法）求解式（7.6.3）比较适宜。理论分析指出，解对称正定方程组的平方根法是一个稳定的算法，在工程计算中使用比较广泛。

追赶法是解三对角方程组（对角元占优）的有效方法，它具有计算量小、方法简单、算法稳定等优点。

关于矩阵的条件数，病态方程组，算法的稳定性，这些概念都是计算数学中比较重要的概念，这里只作了简单的介绍。本章的学习可参看文献
{[}10{]}、{[}11{]}、{[}12{]}。

\begin{quote}
校注（agent 补充，原书疑误）：定理 7.22（2）印作``计算解 \(\overline x\)
精确满足 \((A+\delta A)x=b\)''，此处按原文保留；与上一页同样，方程中的
\(x\) 应按符号约定写为
\(\overline x\)。\href{/Users/xiangjunfeng/Desktop/数值计算专用/numerical-analysis/book/staging/ch07b/assets/211-computed-solution.png}{原式局部图}。

校注（agent
补充，原书疑误）：式（7.6.14）原图在左侧误差比与右侧分式之间未印关系符号，此处保持两式并列；按精度估计用途应为
\(\leqslant\)。另外，由式（7.6.7）得到此界还须
\(\operatorname{cond}(A)_\infty\|\delta A\|_\infty/\|A\|_\infty<1\)，保证分母为正；这一适用条件未在本定理列出。\href{/Users/xiangjunfeng/Desktop/数值计算专用/numerical-analysis/book/staging/ch07b/assets/211-error-estimate.png}{原式局部图}。
\end{quote}

\clearpage

原书 PDF 212；印刷页 199。

\href{/Users/xiangjunfeng/Desktop/数值计算专用/numerical-analysis/book/source-images/pdf-212.jpeg}{原页图像}

\subsection{习题}\label{ux4e60ux9898}

\textbf{1.} 考虑方程组

\[
\begin{cases}
0.409\,6x_1+0.123\,4x_2+0.367\,8x_3+0.294\,3x_4=0.404\,3,\\
0.224\,6x_1+0.387\,2x_2+0.401\,5x_3+0.112\,9x_4=0.155\,0,\\
0.364\,5x_1+0.192\,0x_2+0.378\,1x_3+0.064\,3x_4=0.424\,0,\\
0.178\,4x_1+0.400\,2x_2+0.278\,6x_3+0.392\,7x_4=-0.255\,7.
\end{cases}
\]

（1）用 Gauss 消去法解所给方程组（用四位小数计算）；

（2）用列主元素消去法解所给方程组并且与（1）比较结果。

\textbf{2.} （1）设 \(A\) 是对称阵且 \(a_{11}\neq0\)，经过 Gauss
消去法一步后，\(A\) 约化为
\(\begin{pmatrix}a_{11}&a_1^{\mathrm T}\\0&A_2\end{pmatrix}\)，证明
\(A_2\) 是对称矩阵。

（2）用 Gauss 消去法解对称方程组

\[
\begin{cases}
0.642\,8x_1+0.347\,5x_2-0.846\,8x_3=0.412\,7,\\
0.347\,5x_1+1.842\,3x_2+0.475\,9x_3=1.732\,1,\\
-0.846\,8x_1+0.475\,9x_2+1.214\,7x_3=-0.862\,1.
\end{cases}
\]

\textbf{3.} （1）用式（7.2.9）证明

\[
a_{ij}^{(k)}=a_{ij}^{(1)}-m_{i1}a_{1j}^{(1)}-m_{i2}a_{2j}^{(2)}-\cdots-m_{i,k-1}a_{k-1,j}^{(k-1)}
\quad(i,j\geqslant k),\quad a_{ij}^{(1)}=a_{ij}.
\]

（2）Gauss 消去法中 \(m_{ir},a_{rj}^{(r)}\) 使

\[
m_{ir}=l_{ir}\quad(i>r),\quad a_{rj}^{(r)}=u_{rj}\quad(j\geqslant r),
\]

利用（1）证明

\[
u_{rj}=a_{rj}-\sum_{k=1}^{r-1}l_{rk}u_{kj}\quad(j=r,r+1,\cdots,n),
\]

\[
l_{ir}=\left(a_{ir}-\sum_{k=1}^{r-1}l_{ik}u_{kr}\right)/u_{rr}\quad(i=r+1,\cdots,n).
\]

\textbf{4.} 设 \(A\) 为 \(n\) 阶非奇异矩阵且有分解式 \(A=LU\)，其中
\(L\) 为单位下三角阵，\(U\) 为上三角阵，求证 \(A\)
的所有顺序主子式均不为零。

\textbf{5.} 用 Gauss 消去法说明当
\(\Delta_i\neq0\)（\(i=1,2,\cdots,n-1\)）时，\(A=LU\)，其中 \(L\)
为单位下三角阵，\(U\) 为上三角阵。

\textbf{6.} 设 \(A\) 为 \(n\) 阶矩阵，如果
\(|a_{ii}|>\sum_{\substack{j=1\\j\neq i}}^{n}|a_{ij}|\)（\(i=1,2,\cdots,n\)），则称
\(A\) 为对角占优阵。证明：若 \(A\) 是对角占优阵，经过 Gauss
消去法一步后，\(A\) 具有形式

\[
\begin{pmatrix}a_{11}&a_1^{\mathrm T}\\0&A_2\end{pmatrix},
\]

则 \(A_2\) 是对角占优阵。由此推断，对于对称的对角占优阵来说，用 Gauss
消去法和部分选主元素 Gauss 消去法可得到同样的结果。

\textbf{7.} 设 \(A\) 是对称正定矩阵，经过 Gauss 消去法一步后，\(A\)
约化为
\(\begin{pmatrix}a_{11}&a_1^{\mathrm T}\\0&A_2\end{pmatrix}\)，其中

\[
A=(a_{ij})_n,\quad A_2=(a_{ij}^{(2)})_{n-1}.
\]

证明：（1）\(A\) 的对角元素 \(a_{ii}>0\)（\(i=1,2,\cdots,n\)）；

（第 7 题其余小题续下页。）

\clearpage

原书 PDF 213；印刷页 200。

\href{/Users/xiangjunfeng/Desktop/数值计算专用/numerical-analysis/book/source-images/pdf-213.jpeg}{原页图像}

（续第 7 章习题第 7 题。）

（2）\(A_2\) 是对称正定矩阵；

（3）\(a_{ii}^{(2)}\leqslant a_{ii}\)（\(i=2,3,\cdots,n\)）；

（4）\(A\) 的绝对值最大的元素必在对角线上；

（5）\(\max_{2\leqslant i,j\leqslant n}|a_{ij}^{(2)}|\leqslant\max_{2\leqslant i,j\leqslant n}|a_{ij}|\)；

（6）从（2）、（3）、（5）推出，如果 \(|a_{ij}|<1\)，则对于所有
\(k\)，\(|a_{ij}^{(k)}|<1\)。

\textbf{8.} 设 \(L_k\) 为指标为 \(k\) 的初等下三角阵，即

\[
L_k=\begin{pmatrix}
1&&&&&\\
&\ddots&&&&\\
&&1&&&\\
&&m_{k+1,k}&1&&\\
&&\vdots&&\ddots&\\
&&m_{n,k}&&&1
\end{pmatrix}
\]

（原矩阵上方标注``第 \(k\) 列''，对应 \(m_{k+1,k},\cdots,m_{n,k}\)
所在列。）

（除第 \(k\) 列对角元素之下的元素外，\(L_k\) 和单位阵 \(I\)
相同），求证：当 \(i,j>k\) 时，\(\widetilde L_k=I_{ij}L_kI_{ij}\)
也是一个指标为 \(k\) 的初等下三角阵，其中 \(I_{ij}\) 为初等排列阵。

\textbf{9.} 试推导矩阵 \(A\) 的 Crout 分解 \(A=LU\) 的计算公式，其中
\(L\) 为下三角阵，\(U\) 为单位上三角阵。

\textbf{10.} 设 \(Ux=d\)，其中 \(U\) 为三角阵。

（1）就 \(U\) 为上及下三角阵推导一般的求解公式，并写出算法。

（2）计算解三角方程组 \(Ux=d\) 的乘除法次数。

（3）设 \(U\) 为非奇异阵，试推导求 \(U^{-1}\) 的计算公式。

\textbf{11.} 证明：（1）如果 \(A\) 是对称正定矩阵，则 \(A^{-1}\)
也是对称正定矩阵；

（2）如果 \(A\) 是对称正定矩阵，则 \(A\) 可唯一地写成
\(A=L^{\mathrm T}L\)，其中 \(L\) 是具有正对角元的下三角阵。

\textbf{12.} 用 Gauss-Jordan 方法求 \(A\) 的逆阵，其中

\[
A=\begin{pmatrix}
2&1&-3&-1\\
3&1&0&7\\
-1&2&4&-2\\
1&0&-1&5
\end{pmatrix}.
\]

\textbf{13.} 用追赶法解三对角方程组 \(Ax=b\)，其中

\[
A=\begin{pmatrix}
2&-1&0&0&0\\
-1&2&-1&0&0\\
0&-1&2&-1&0\\
0&0&-1&2&-1\\
0&0&0&-1&2
\end{pmatrix},\quad
b=\begin{pmatrix}1\\0\\0\\0\\0\end{pmatrix}.
\]

\textbf{14.} 用改进的平方根法解方程组

\[
\begin{pmatrix}
2&-1&1\\
-1&-2&3\\
1&3&1
\end{pmatrix}
\begin{pmatrix}x_1\\x_2\\x_3\end{pmatrix}
=\begin{pmatrix}4\\5\\6\end{pmatrix}.
\]

\textbf{15.} 下述矩阵能否分解为 \(LU\)（其中 \(L\) 为单位下三角阵，\(U\)
为上三角阵）？若能分解，那么分解是否唯一？

\[
A=\begin{pmatrix}1&2&3\\2&4&1\\4&6&7\end{pmatrix},\quad
B=\begin{pmatrix}1&1&1\\2&2&1\\3&3&1\end{pmatrix},\quad
C=\begin{pmatrix}1&2&6\\2&5&15\\6&15&46\end{pmatrix}.
\]

\clearpage

原书 PDF 214；印刷页 201。

\href{/Users/xiangjunfeng/Desktop/数值计算专用/numerical-analysis/book/source-images/pdf-214.jpeg}{原页图像}

（续第 7 章习题。）

\textbf{16.} 试划出部分选主元素三角分解法框图，并且用此法解方程组

\[
\begin{pmatrix}0&3&4\\1&-1&1\\2&1&2\end{pmatrix}
\begin{pmatrix}x_1\\x_2\\x_3\end{pmatrix}
=\begin{pmatrix}1\\2\\3\end{pmatrix}.
\]

\textbf{17.} 如果方阵 \(A\) 有 \(a_{ij}=0\)（\(|i-j|>t\)），则称 \(A\)
为带宽 \(2t+1\) 的带状矩阵。设 \(A\) 满足三角分解条件，试推导 \(A=LU\)
的以下计算公式：

对于 \(r=1,2,\cdots,n\)，有

（1）

\[
u_{ri}=a_{ri}-\sum_{k=\max(1,i-t)}^{r-1}l_{rk}u_{ki}\quad(i=r,r+1,\cdots,\min(n,r+t));
\]

（2）

\[
l_{ir}=\left(a_{ir}-\sum_{k=\max(1,i-t)}^{r-1}l_{ik}u_{kr}\right)/u_{rr}\quad(i=r+1,\cdots,\min(n,r+t)).
\]

\textbf{18.} 设 \(A=\begin{pmatrix}0.6&0.5\\0.1&0.3\end{pmatrix}\)，计算
\(A\) 的行范数、列范数、2-范数及 \(F\)-范数。

\textbf{19.}
求证：（1）\(\|x\|_\infty\leqslant\|x\|_1\leqslant n\|x\|_\infty\)；（2）\(\frac1{\sqrt n}\|A\|_F\leqslant\|A\|_2\leqslant\|A\|_F\)。

\textbf{20.} 设 \(P\in\mathbf R^{n\times n}\) 且非奇异，又设 \(\|x\|\)
为 \(\mathbf R^n\) 上一向量范数，定义 \(\|x\|_P=\|Px\|\)。试证明
\(\|x\|_P\) 是 \(\mathbf R^n\) 上向量的一种范数。

\textbf{21.} 设 \(A\in\mathbf R^{n\times n}\) 为对称正定矩阵，定义
\(\|x\|_A=(Ax,x)^{\frac12}\)，试证明 \(\|x\|_A\) 为 \(\mathbf R^n\)
上向量的一种范数。

\textbf{22.} 证明：当且仅当 \(x\) 和 \(y\) 线性相关且
\(x^{\mathrm T}y\geqslant0\) 时，才有 \(\|x+y\|_2=\|x\|_2+\|y\|_2\)。

\textbf{23.} 分别描述 \(\mathbf R^2\)
中（画图）\(S_v=\{x\mid\|x\|_v=1,x\in\mathbf R^2\}\)（\(v=1,2,\infty\)）。

\textbf{24.} 令 \(\|\cdot\|\) 是 \(\mathbf R^n\)（或
\(\mathbf C^n\)）上的任意一种范数，而 \(P\)
是任一非奇异实（或复）矩阵，定义范数 \(\|x\|'=\|Px\|\)，证明
\(\|A\|'=\|PAP^{-1}\|\)。

\textbf{25.} 设 \(\|A\|_s,\|A\|_t\) 为 \(\mathbf R^{n\times n}\)
上任意两种矩阵算子范数，证明存在常数 \(c_1,c_2>0\)，使对于一切
\(A\in\mathbf R^{n\times n}\) 满足

\[
c_1\|A\|_s\leqslant\|A\|_t\leqslant c_2\|A\|_s.
\]

\textbf{26.} 设 \(A\in\mathbf R^{n\times n}\)，求证 \(A^{\mathrm T}A\)
与 \(AA^{\mathrm T}\) 特征值相等，即求证
\(\lambda(A^{\mathrm T}A)=\lambda(AA^{\mathrm T})\)。

\textbf{27.} 设 \(A\) 为非奇异矩阵，求证

\[
\frac1{\|A^{-1}\|_\infty}=\min_{y\neq0}\frac{\|A\|_{y\infty}}{\|y\|_\infty}.
\]

\textbf{28.} 设 \(A\) 为非奇异矩阵，且
\(\|A^{-1}\|\|\delta A\|<1\)，求证 \((A+\delta A)^{-1}\) 存在且有估计

\[
\frac{\|A^{-1}-(A+\delta A)^{-1}\|}{\|A^{-1}\|}
\leqslant
\frac{\operatorname{cond}(A)\dfrac{\|\delta A\|}{\|A\|}}
{1-\operatorname{cond}(A)\dfrac{\|\delta A\|}{\|A\|}}.
\]

\textbf{29.} 矩阵第 1 行乘以一个数，成为
\(A=\begin{pmatrix}2\lambda&\lambda\\1&1\end{pmatrix}\)，证明当
\(\lambda=\pm\frac23\) 时，\(\operatorname{cond}(A)_\infty\) 有最小值。

\textbf{30.} 设 \(A\) 为对称正定矩阵，且其分解为
\(A=LDL^{\mathrm T}=W^{\mathrm T}W\)，其中
\(W=D^{1/2}L^{\mathrm T}\)，求证：

（1）\(\operatorname{cond}(A)_2=[\operatorname{cond}(W)_2]^2\)；（2）\(\operatorname{cond}(A_2)=\operatorname{cond}(W^{\mathrm T})_2\operatorname{cond}(W)_2\)。

\textbf{31.} 设 \(A=\begin{pmatrix}100&99\\99&98\end{pmatrix}\)，计算
\(A\) 的条件数 \(\operatorname{cond}(A)_v\)（\(v=2,\infty\)）。

\textbf{32.} 证明：如果 \(A\) 是正交矩阵，则
\(\operatorname{cond}(A)_2=1\)。

\textbf{33.} 设 \(A,B\in\mathbf R^{n\times n}\) 且 \(\|\cdot\|\) 为
\(\mathbf R^{n\times n}\) 上矩阵的算子范数，证明
\(\operatorname{cond}(AB)\leqslant\operatorname{cond}(A)\operatorname{cond}(B)\)。

\begin{quote}
校注（agent 补充，原书疑误）：第 27 题分子中的 \(y\)
在原图排到闭合范数竖线外的下标位置，与 \(\infty\) 并列，此处以
\(\|A\|_{y\infty}\) 保留这一印刷位置。按题意，该分子应为
\(\|Ay\|_\infty\)。\href{/Users/xiangjunfeng/Desktop/数值计算专用/numerical-analysis/book/staging/ch07b/assets/214-ex27.png}{原式局部图}。

校注（agent 补充，原书疑误）：第 30 题（2）左端原图印为
\(\operatorname{cond}(A_2)\)，此处保留；同题只定义了矩阵 \(A\)，并使用
2-范数条件数，左端应为
\(\operatorname{cond}(A)_2\)。\href{/Users/xiangjunfeng/Desktop/数值计算专用/numerical-analysis/book/staging/ch07b/assets/214-ex30.png}{原式局部图}。
\end{quote}

\end{document}
